\section{从 AlphaStar 到 LLM：协作用户与对抗用户共存}

League 机制来自对抗游戏，但 LLM 部署并非纯零和。本节把 population 拆成合作用户、攻击者、熟悉任务和新任务，讨论如何在帮助正常用户与限制恶意利用之间建立联合目标，而不是把所有交互都强行解释成“赢或输”。

在 54 分钟后，讲者转入直接对照：LLM 环境不是标准零和博弈，用户大多是协作者，但对抗用户同样真实存在。对应地，训练系统必须同时处理两类目标。

\begin{figure}[H]
\centering
\includegraphics[width=0.92\textwidth]{frame-05.jpg}
\caption{讲者对当前 LLM 场景的对照总结}
\label{fig:contrast-llm}
\end{figure}

\begin{knowledgebox}{读图：LLM 对照页给出三项迁移警告}
首先，真实奖励无法像游戏胜负那样完整指定；其次，LLM 计算预算长期偏向 pretraining；最后，任务与用户分布远比标准地图开放。多 Agent red teaming 因此只是训练组织方式之一，不能替代产品目标、权限边界和真实结果评测。
\end{knowledgebox}

\begin{figure}[H]
\centering
\includegraphics[width=0.95\textwidth]{diagrams/user-distribution.png}
\caption{讲义整理图：部署 population 中的合作/对抗用户与熟悉/新颖任务}
\label{fig:user-distribution}
\end{figure}

\begin{knowledgebox}{读图：四个象限必须分别报告}
合作用户上的高成功率不能证明系统抗攻击，熟悉任务上的鲁棒性也不能证明能处理新工具。联合目标应同时包含 helpfulness、robustness、cost 和权限约束，并报告各子群体最差表现。平均分会掩盖“主流用户很好、少数高风险场景失守”的部署缺陷。
\end{knowledgebox}
图 \ref{fig:contrast-llm} 的要点非常清晰\footnote{来源：\href{https://www.youtube.com/watch?v=ntjOxjZMaac}{Lecture 03 视频帧}，时间戳 00:56:30。}：用户通常合作，但必须假设 adversarial probing 常态化；现实任务复杂度远高于游戏；多 Agent red teaming 已在使用，但整体仍处于早期。

\begin{knowledgebox}{迁移时最容易忽视的差异}
\begin{itemize}
    \item 奖励函数模糊：真实任务很少有明确胜负信号。
    \item Prompt 空间无界：自然语言比离散动作接口复杂得多。
    \item 用户目标异质：同一模型同时服务效率、创造力、安全、合规等多目标。
\end{itemize}
\end{knowledgebox}

\begin{importantbox}{双生态训练目标}
迁移后的 LLM 多 Agent 训练至少应同时优化两类分布：
\begin{enumerate}
    \item \textbf{Collaborative distribution}：正常用户任务成功率、成本和体验。
    \item \textbf{Adversarial distribution}：越权、诱导、注入、策略欺骗下的稳健性。
\end{enumerate}
单看一类分布会导致另一类指标崩塌。
\end{importantbox}

\begin{warningbox}{错误策略：把 red teaming 当作发布前一次性流程}
讲者强调 ``users probe it''，意味着攻击会持续演化。red teaming 必须从项目流程角度变成持续对抗循环，而不是 QA 阶段的 checklist。
\end{warningbox}

这里的关键引语也很值得保留：\textit{``Users are in general collaborative... That being said, we also have adversarial users.''}\footnote{来源：\href{https://www.youtube.com/watch?v=ntjOxjZMaac}{Lecture 03 字幕}，区间 00:54:47--00:55:12。}

\subsection{本章小结}
LLM 时代的多 Agent 问题不是把所有人当敌人，也不是把所有人当朋友，而是承认二者共存。训练、评测与上线监控都必须对应这一现实分布。

